\section{The books and the calendar}

\subsection{Finding the day}

This handbook follows the postconciliar Roman Rite in the dioceses of the United States. The books for Mass are the \emph{Missale Romanum}, third typical edition (2002, emended 2008), its U.S. English edition (2011), and the U.S. \emph{Lectionary for Mass}, second typical edition (Sunday volume, 1998, weekday and other volumes, 2001). For the Hours use the Proper of Seasons in volume I of the U.S. \emph{Liturgy of the Hours}. The Missal supplies the priest's prayers. The Lectionary appoints the passages proclaimed from the ambo. In the Hours the Church prays at appointed times of day. The reading lists use the U.S. Lectionary’s book names and verse divisions. Older Catholic Bibles can number a psalm or divide a prophetic book differently. The home prayers give their Douay references. With a Bible and these references you can keep the home course even if you have none of the liturgical books.

Advent opens with Evening Prayer I of the Sunday nearest November 30, and ends before Evening Prayer I of the Nativity. The first Sunday therefore falls between November 27 and December 3. On a Sunday use the letter in the parish missal. To find it by reckoning, take the calendar year of the Christmas you are approaching and add one. Divide that number by three. A remainder of one gives Year A, two gives Year B, and zero gives Year C. The weekday readings of Advent recur annually. The I and II beside weekday readings in Ordinary Time have no bearing on this course. A celebration proper to your church or diocese can determine what you hear there.

Suppose December 17 falls on Sunday. At Mass you hear the Third Sunday of Advent, and at Evening Prayer you can listen for the first O Antiphon. On Monday morning turn to December 18. A week later December 24 is the Fourth Sunday. Go to the Sunday readings in the morning and to the Christmas Vigil in the evening. In a longer Advent, the Fourth Sunday can fall on December 18. On Monday turn directly to December 19. The Church gives us the days that occur in this Advent. There is no need to make up a missing weekday.

\subsection{Praying with someone who uses the 1962 Missal}

Friends who attend the two Masses can pray Mary's canticle together on the same evening. The next morning they may hear quite different parts of the Gospel. Take the First Sunday. In the 1962 Missal, worshippers hear Paul call the Romans to wake from sleep and Luke recount the coming of the Son of Man (Romans 13:11--14a and Luke 21:25--33). In the present Lectionary the Roman epistle is heard on that Sunday in Year A, with Matthew. Luke is heard in Year C, with Jeremiah and a letter to the Thessalonians. In Year B, Mark sets a doorkeeper at his post. Reading the whole passage appointed in either book allows the evangelist's particular counsel to be heard.

In the older Missal an ordinary Advent feria, a weekday without a proper Mass or other celebration taking precedence, generally returns to the Mass of the previous Sunday. Wednesday, Friday and Saturday after the Third Sunday are the December Ember Days, with their own lessons. On Ember Wednesday the promise to Ahaz and the angel's visit to Mary are heard together. In the present U.S. Lectionary, Isaiah 7:10--14 and Luke 1:26--38 meet on the weekday of December 20. For the older date find the Third Sunday and count forward. For the present appointment find December 20 and check whether it is Sunday. And on a Sunday morning that is December 24 the 1962 books give the Vigil of the Nativity, while the present books give the Fourth Sunday. One set of directions could easily send a family to the wrong Gospel on the very morning they most wanted to prepare.

The 1962 Breviary places the lessons of Matins within its own sequence of psalms and nocturns. In the revised Hours, a reader can turn to the Office of Readings for the biblical and patristic lessons and pray it at a suitable hour of the day. Both books give us Isaiah, the Baptist and Our Lady. Both have put on Christian lips the plea for the coming of Emmanuel. A Catholic visiting a different parish can therefore bring the same faith, hopes and works of mercy, and open the volume that follows the books on the altar and at the ambo that day.

\subsection{Friday penance in the United States}

The universal discipline for Latin Catholics is found in the 1983 \emph{Code of Canon Law}, canons 1249--1253, promulgated by St John Paul II. Canon 1251 prescribes Friday abstinence unless a solemnity occurs, and fast and abstinence on Ash Wednesday and Good Friday. Canon 1253 entrusts the determination of further particulars and substitutions to episcopal conferences. In its \emph{Pastoral Statement on Penance and Abstinence} (1966), nos.~7, 22--25, the U.S. conference urged prayer and works of charity in Advent, retained Friday as a day of penance, and gave pride of place to voluntary abstinence from meat on ordinary Fridays outside Lent. See also the U.S. complementary norm for canons 1252--1253. The account of discipline here was checked on November 23, 2026. The dates in the historical chapters identify the witnesses to earlier practice.
